% Generated by report_gen.py (hwde v1) - do not hand-edit.
\documentclass[11pt]{article}
\newcommand\hspaper{a4paper}
\usepackage[nodiagramkit]{housestyle}
\usepackage{longtable}
\usepackage{pdfpages}
\makeatletter
\@ifundefined{KV@Gin@artifact}{\define@key{Gin}{artifact}[]{}}{}
\makeatother
\setcounter{tocdepth}{1}
\hsslug{PCB-0020-A \textperiodcentered\ esp32c3-node}
\renewcommand{\sectionmark}[1]{\markboth{#1}{}}
\hssection{\leftmark}
\begin{document}
\sloppy
\hstitleblock{PCB-0020-A \textperiodcentered\ hwde v1 pipeline \textperiodcentered\ generated 2026-10-01 09:03}%
{esp32c3-node \textemdash{} Highlights}%
{Phase P8; all 6 recorded gates pass.}


\section{At a Glance}
A tiny WiFi + Bluetooth LE sensor node, in the spirit of the Seeed XIAO ESP32-C3. Product scope (no mode token), 2 layers, meant to go to JLCPCB as-is with economy PCBA.

\begin{itemize}
\item ESP32-C3-MINI-1 module; keep its antenna keep-out clear of copper and parts, antenna at a board edge.
\item USB-C for power and flashing over the C3's native USB (no USB-UART bridge); BOOT and RESET buttons.
\item Runs from a single-cell LiPo (JST-PH 2-pin) with a charger that has power path, so it runs
from USB with or without a battery. 3.3 V LDO for everything. Battery voltage sense into an ADC pin.
\item A BME280-class temperature/humidity/pressure sensor on I2C, plus a Qwiic/STEMMA QT connector on the same bus.
\item One status LED on a GPIO.
\item Keep it small, roughly XIAO-sized if an honest layout allows; spare GPIOs on castellated or 0.1 in edge pads.
\end{itemize}
{\small
\begin{longtable}{lp{11cm}}
\hstoprule
\hshead{Field} & \hshead{Value} \\
\hline
\endhead
part number & PCB-0020-A \\
board & esp32c3-node \\
phase & P8 (Verification) \\
size & 26.0 x 38.0 mm \\
gates & 6 of 6 recorded gates pass \\
\hline
\end{longtable}
}
\section{The Board}
\begin{center}
\includegraphics[width=0.48\textwidth]{reports/renders/esp32c3-node_top.png}\hfill
\includegraphics[width=0.48\textwidth]{reports/renders/esp32c3-node_bottom.png}\hfill
\\{\small top and bottom, as rendered from the board file}
\end{center}
\section{Copper Layers}
The outer copper layers as KiCad shows them, with every pad, track and zone labelled with its net (zoom in to read the small ones).
\clearpage
\subsection*{F.Cu}
\begin{center}
\includegraphics[width=\textwidth,height=0.82\textheight,keepaspectratio]{reports/layers/PCB-0020-A_esp32c3-node_F_Cu.pdf}\\
{\small\texttt{reports/layers/PCB-0020-A\_esp32c3-node\_F\_Cu.pdf}}
\end{center}
\clearpage
\subsection*{B.Cu}
\begin{center}
\includegraphics[width=\textwidth,height=0.82\textheight,keepaspectratio]{reports/layers/PCB-0020-A_esp32c3-node_B_Cu.pdf}\\
{\small\texttt{reports/layers/PCB-0020-A\_esp32c3-node\_B\_Cu.pdf}}
\end{center}
\subsection*{3D view}
\begin{center}
\includegraphics[width=0.8\textwidth]{reports/layers/PCB-0020-A_esp32c3-node_iso.png}\\
{\small\texttt{reports/layers/PCB-0020-A\_esp32c3-node\_iso.png}}
\end{center}
\section{What Is On It}
\includepdf[pages=1]{reports/schematic.pdf}
\section{Key Decisions}
{\small
\begin{longtable}{lp{1.6cm}p{6.2cm}p{6.2cm}}
\hstoprule
\hshead{Ph} & \hshead{When} & \hshead{Decision} & \hshead{Why} \\
\hline
\endhead
P0 & 10-01 05:36 & Q1b cell class \textasciitilde{}400 mAh (300-1000 mAh acceptable); PROVISIONAL (delegate, owner said decide for yourself; re-confirm at H1) & XIAO-sized node class; sizing the charge current to the smallest plausible cell is the conservative choice \\
P2 & 10-01 05:48 & accepted deviation: VBUS attach charge \textasciitilde{}65 uC \textgreater{} USB2.0 50 uC guideline & LDO current limit bounds surge; shown at H1 \\
P3 & 10-01 06:11 & fp\_verify U1 (C2838502) exit 1 'pad count 61 != 53' accepted as a checker artifact, not a footprint defect & Footprint has exactly 53 pad numbers; pad 49 (GND EPAD) is a 3x3 array of 1.45 mm squares (one custom pad chamfered for pin-1 orientation, reported as 0.005 mm), total 5.4 mm square as DS Fig 11-1. Edge pads 0.4x0.8 at 0.8 pitch, corners 0.7x0.7, all match.... \\
P3 & 10-01 06:11 & J1 C165948, J2 C160404, J3 C295747 footprints accepted from the JLC/EasyEDA pull without a datasheet extraction & No prior board has extractions for them; spend-bounded run. Pad geometry checked by hand: J2 1.0 pitch 0.6x1.55 + 1.2x1.8 mounts; J3 2.0 pitch 1.0x3.8 + 1.5x3.4 mounts; J1 16-pin USB-C 0.3x1.3 at 0.5 pitch + 4 shell THT. Re-check against the vendor drawings... \\
\hline
\end{longtable}
}
\section{How the Run Went}
23 recorded decisions from P0 to P3, 2026-10-01T05:36 to 2026-10-01T06:11. The run went back to an earlier phase 0 times.
\section{Checks and Cost}
{\small
\begin{longtable}{llrr}
\hstoprule
\hshead{Gate} & \hshead{Status} & \hshead{Attempts} & \hshead{Failing} \\
\hline
\endhead
erc & pass & 5 & 0 \\
place & pass & 5 & 0 \\
drc\_routed & pass & 11 & 0 \\
verify & pass & 7 & 0 \\
dfm & pass & 2 & 0 \\
\hline
\end{longtable}
}
Model cost of generating this board: \$24.98.

The full design doc (esp32c3-node-design-full) has every decision, the renders and the run's history.

\end{document}
